% !TeX root = ../main.tex
% Chapter 7 -- Expected Outcomes, Gate-Conditioned Decisions, and Risk Analysis
% Self-contained section file: no preamble, no \documentclass.
% HONESTY RULE: no experiments have been run. This chapter replaces a "Results"
% chapter with an EXPECTED / GATE-CONDITIONED OUTCOMES framing. Every quantitative
% figure that appears as a number belongs to the LITERATURE and is attributed as such.
%
% PREAMBLE REQUIREMENTS (main.tex must load these):
%   \usepackage{booktabs}
%   \usepackage{tikz}
%   \usetikzlibrary{arrows.meta,positioning}
% The TikZ gate tree uses arrows.meta (Latex tips) and positioning (below=of / right=of).

\chapter{Expected Outcomes, Gate-Conditioned Decisions, and Risk Analysis}
\label{ch:expected-outcomes}

\newcommand{\evangelistanums}{AUC~$0.94$ / sensitivity~$90.7\%$ / specificity~$86.6\%$}

\begin{quote}\small\itshape
No experiments have been run, and the no-results stance is structural rather than a
stylistic choice: the corpus has not been forked (the repository's \texttt{data/}
substrate is empty), no independent per-AU veterinary anchor exists, and the only
available source is a Roboflow \emph{binary} pain corpus with no native $0$--$10$ FGS
labels. This chapter therefore contains no measured results, no accuracies, no
$\kappa$ values, no calibration numbers, and no figures of our own curves. It states,
per gate, the \emph{decision rule} and the gate-conditioned outcome
(\textsc{go} / \textsc{pivot} / \textsc{collapse}), the modal product we plan to ship,
and the kill criteria under which we pivot. Every numeric value that appears
(e.g. \evangelistanums) is drawn from the prior literature and is explicitly
attributed to its authors as a \emph{background anchor}, never as a result of this
work. Real-cat evaluation is named throughout as future work; the present
contribution is a methods/protocol specification.
\end{quote}

This chapter deliberately departs from the template's Chapter~4
(\emph{การทดลองและผลลัพธ์}, Experiments and Results). Because the project is a
research proposal and methodology with a strict, artifact-emitting gate order
rather than a completed empirical study, a conventional results chapter would
either be empty or be forced to fabricate numbers. We therefore recast the
chapter as an \emph{expected-outcomes and risk} analysis: for each gate we
pre-register the decision rule, the pass/fail branch, and exactly what ships in
each case. The chapter is organised around the \emph{modal deliverable}
(\S\ref{sec:modal-deliverable}), the gate-conditioned decision table
(\S\ref{sec:decision-table}), the power-calculation-conditioned framing of the
abstention curve (\S\ref{sec:power-conditioned}), the expected severity-cell
collapse (\S\ref{sec:severity-collapse}), the Monte-Carlo correlation-sensitivity
band (\S\ref{sec:rho-band}), and the literature anchors we cite as background
(\S\ref{sec:lit-anchors}).

%====================================================================
\section{The Modal Deliverable as the Default Product}
\label{sec:modal-deliverable}

We plan the repository, the evaluation protocol, and the write-up for the
\emph{modal} outcome --- the product we expect to ship --- not the optimistic one.
The modal deliverable is a \textbf{binary-pain decision-support instrument with a
trustworthiness wrapper}, headlined by a power-conditioned, per-AU
confound-attribution protocol and accompanied by a guarded VLM-as-AU-rater $\kappa$
reliability check, shipped on an explicitly-conceded engine. The graded $0$--$10$ Feline Grimace Scale (FGS) layer
is present strictly as an \emph{inspected-not-validated} artifact. The single
governing rule on the prose: \textbf{the abstract must hold with the graded layer
dropped}, and the word ``graded'' is struck from every validated-claim sentence.

Concretely, the modal product comprises six named components, summarised in
Table~\ref{tab:modal-components}. Four are validated claims (Components~1--4);
Component~5 (the NPV/abstention curve) is reported only as placeholder-illustrative
sizing, never as a validated number; and the sixth (the graded CORN layer) is never
entered into a validated results table.

\begin{table}[t]
\centering
\caption[The six components of the modal deliverable]{The six components of the
modal deliverable. The abstract is required to hold with Component~6 dropped; the
DINOv2+CORN engine is conceded as plumbing and is never claimed novel. ``Validated''
means estimated against vet-confirmed labels under the circularity firewall.}
\label{tab:modal-components}
\small
\begin{tabular}{@{}clp{0.40\linewidth}c@{}}
\toprule
\# & Component & What it is & Validated? \\
\midrule
1 & Calibrated binary decision &
  Per-image pain decision at a \emph{fixed pain-recall} $\geq 0.90$ operating
  point (after \citealp{evangelista2019}); never Youden-$J$/$F_1$ & Yes \\
2 & Confound-attribution protocol &
  FGS-BG-Gap (background-swap primitive \citep{xiao2021noise,moayeri2022comprehensive})
  $+$ per-AU EBPG (Score-CAM energy-based pointing-game primitive
  \citep{wang2020scorecam}; distinct from the per-AU saliency-alignment framing of
  \citet{lencioni2025segment}), \emph{one-directional} (``no confound detected at
  this power'' \citep{adebayo2022posthoc}) & Yes (headline) \\
3 & $\kappa$ reliability check &
  Per-AU VLM-vs-vet quadratic-weighted $\kappa$ \citep{cohen1968kappa} reported on
  its \emph{CI lower bound} --- a standard clinimetric acceptance gate
  \citep{tractenberg2010kappa,donner2010kappa,rotondi2012kappa,rotondi2018kappasize,simwright2005kappa},
  not novel here; guarded, inspected-not-validated & Yes (guarded check) \\
4 & Welfare-asymmetric decision curve &
  Net-benefit decision-curve analysis \citep{vickers2006dca} with the
  undertreat:overtreat harm ratio swept as a \emph{range} & Yes \\
5 & LB-abstention curve &
  One-sided $95\%$ NPV \emph{lower bound} at each abstention rate
  \citep{angelopoulos2021ltt}; ``guaranteed'' is banned & Illustrative (placeholder sizing) \\
6 & Inspected graded-CORN &
  $5$ per-AU CORN heads \citep{shi2023corn} $\to$ $0$--$10$ sum $\to$ pmf;
  inspected-not-validated & \textbf{No} \\
\bottomrule
\end{tabular}
\end{table}

The headline contribution is Component~2 (the power-conditioned, per-AU
confound-attribution protocol); Component~3 (the VLM-as-AU-rater $\kappa$) is a
guarded reliability check held below it as kill-tree insurance, not a co-equal
second pillar. Both ship as released, runnable, dataset-agnostic code. Novelty is
claimed for the \emph{assembly} and the per-AU EBPG-as-confound-evidence
instantiation of equivalence-style audit hygiene, never for any individual primitive
or the one-directional / power-conditioned statistics themselves
\citep{huanghooker2026fairness,singh2023samplesize}. The calibration, RPS,
ClasswiseECE, and conformal machinery are \emph{cited supporting evidence}, never
headlines. The engine (frozen DINOv2 ViT-S/14 --- with the \texttt{\_reg}
register variant (\texttt{dinov2\_vits14\_reg}) the field default, A/B'd before
locking ViT-S/14 because registers suppress attention artifacts that hurt the
dense, localized per-AU features
\citep{oquab2023dinov2} feeding the CORN ordinal heads \citep{shi2023corn}) is
conceded as plumbing and is never presented as novel.

%====================================================================
\section{The Gate-Conditioned Decision Table}
\label{sec:decision-table}

The project executes a strict, blocking gate order: \textbf{G0}
(power calcs and vet-budget pre-registration) $\to$ \textbf{G1} (per-CAT merge)
$\to$ \textbf{G2} (capture-condition confound audit) $\to$ \textbf{G3} (frozen
hashed cat-disjoint hold-out) $\to$ \textbf{G4} (MPS compute-correctness, blocking)
$\to$ \textbf{G5} (alignment/NME) $\to$ \textbf{G1-B} (weak-label reliability
$\kappa$ pilot) $\to$ \textbf{G6} (severity-cell collapse). Each gate carries a
\emph{pre-registered} decision rule; passing a gate is a precondition for the work
downstream of it, and failing a gate routes to a pre-committed branch rather than
to an improvised fix. Figure~\ref{fig:gate-tree} renders the gate order as a
decision tree terminating in the modal product, and Table~\ref{tab:kill-crosswalk}
gives the kill/pivot crosswalk in full.

\begin{figure}[t]
\centering
\begin{tikzpicture}[
  node distance=6.5mm and 9mm,
  every node/.style={font=\footnotesize},
  gate/.style={rectangle, rounded corners, draw, fill=blue!6, align=center,
               minimum width=30mm, minimum height=8mm, inner sep=3pt},
  branch/.style={rectangle, draw, fill=orange!8, align=center,
                 minimum width=34mm, inner sep=3pt},
  ship/.style={rectangle, rounded corners, draw, very thick, fill=green!10,
               align=center, minimum width=42mm, inner sep=4pt},
  flow/.style={-{Latex[length=2mm]}, thick},
  sidelab/.style={font=\scriptsize\itshape}
]
% main spine
\node[gate] (g0) {G0 power / vet budget};
\node[gate, below=of g0] (g1) {G1 per-CAT merge};
\node[gate, below=of g1] (g2) {G2 confound audit};
\node[gate, below=of g2] (g3) {G3 frozen hold-out};
\node[gate, below=of g3] (g4) {G4 MPS correctness\\\scriptsize(BLOCKING)};
\node[gate, below=of g4] (g5) {G5 alignment / NME\\\scriptsize(BLOCKING)};
\node[gate, below=of g5] (g1b) {G1-B $\kappa$ pilot\\\scriptsize(CI lower bound)};
\node[gate, below=of g1b] (g6) {G6 severity cells};
\node[ship, below=of g6] (modal) {Modal product:\\binary $+$ wrapper\\$+$ confound protocol (headline)\\$+$ guarded $\kappa$ check\\$+$ graded (inspected-only)};

\foreach \a/\b in {g0/g1, g1/g2, g2/g3, g3/g4, g4/g5, g5/g1b, g1b/g6, g6/modal}
  \draw[flow] (\a) -- (\b);

% branch outcomes on the right
\node[branch, right=18mm of g0] (g0b) {G0(b):\\NPV band ships\\placeholder-illustrative};
\node[branch, right=18mm of g2] (g2b) {G2 non-portable:\\demote to\\threat-to-validity};
\node[branch, right=18mm of g4] (g4b) {G4/G5 fail:\\\textsc{collapse} --\\fix, believe no number};
\node[branch, right=18mm of g1b] (g1bb) {G1-B $\kappa$-LB low:\\\textsc{pivot} to\\binary spine};
\node[branch, right=18mm of g6] (g6b) {G6 fires:\\\textsc{collapse}\\high-end scale};

\draw[flow, orange!70!black] (g0) -- (g0b);
\draw[flow, orange!70!black] (g2) -- (g2b);
\draw[flow, orange!70!black] (g4) -- (g4b);
\draw[flow, orange!70!black] (g1b) -- (g1bb);
\draw[flow, orange!70!black] (g6) -- (g6b);

% pivots still terminate in a publishable product
\draw[flow, dashed, orange!70!black] (g1bb.south) |- ([yshift=4mm]modal.east) -- (modal.east);
\end{tikzpicture}
\caption[Gate-conditioned decision tree]{Gate-conditioned decision tree. The left
spine is the \textsc{go} path through gates G0$\to$G6 terminating in the modal
product. The right column shows the pre-registered branch at each gate. Note that
no branch is a project-killing null: a G1-B \textsc{pivot} drops the graded layer
but still ships the calibrated binary spine, the headline confound protocol, and
the guarded $\kappa$ check (dashed arrow). G4/G5 are blocking correctness gates --- a
failure halts reporting until fixed, never silently degrades a number.}
\label{fig:gate-tree}
\end{figure}

\begin{table}[t]
\centering
\caption[Kill/pivot crosswalk]{Kill/pivot crosswalk (gate $\to$ branch $\to$ what
ships), after the kill criteria of the finalised direction. Every branch remains
publishable; the only ``hard stop'' branches are the blocking correctness gates
G4/G5, which halt reporting until the defect is fixed rather than ending the
project.}
\label{tab:kill-crosswalk}
\footnotesize
\begin{tabular}{@{}p{0.13\linewidth}p{0.36\linewidth}p{0.42\linewidth}@{}}
\toprule
Gate & Failing branch (trigger) & What ships on that branch \\
\midrule
G0(b) &
  NPV sizing rests on a placeholder budget over zero data (not yet certifiable) &
  \textsc{proceed}; the abstention curve (Component~5) is reported as
  \emph{placeholder-illustrative} with a one-sided $95\%$ lower bound;
  ``guaranteed'' never appears. Not a kill --- relocates the deliverable onto the
  welfare decision curve. \\[2pt]
G2 &
  Audit damning but yields only ``CAT\_01 is confounded'' (non-portable) &
  Demote the audit to a threat-to-validity paragraph; the guarded $\kappa$ check
  becomes the sole-survivor headline (kill-tree insurance). If G1-B \emph{also}
  fails, the honest paper is a confound/audit note (venue pre-registered now). \\[2pt]
G2 &
  Confounding detected as a \emph{transportable protocol} &
  \textsc{proceed} --- this is the headline leg, publishable even if everything
  downstream is null \citep{xiao2021noise,moayeri2022comprehensive}. \\[2pt]
G1-B &
  $\kappa$ CI lower bound below floor on orbital/ear/head &
  \textsc{pivot} to binary spine: \emph{drop graded entirely}; ship calibrated
  binary $+$ abstention $+$ the headline confound protocol. The guarded $\kappa$
  check and the cited wrapper still stand (the $\kappa$ protocol remains runnable,
  with its result pending the independent vet anchor). \\[2pt]
G1-B &
  Passes orbital/ear/head; fails muzzle/whiskers &
  \textsc{proceed} on the surviving AUs; state how dropping $2$ of $5$ heads
  shifts the achievable $0$--$10$ range and whether $0.39$ ($\approx 4/10$) is
  even reachable. \\[2pt]
G4 / G5 &
  MPS logit parity / decode unit test / NME audit fails &
  \textsc{blocking}: fix before any metric is believed. No number from
  silently-wrong logits or misaligned crops is reported. \\[2pt]
G6 &
  AU=2 (severe) cells single-digit &
  \textsc{collapse} the high-end scale (pre-committed); do \emph{not} print a
  per-cell severe calibration number. Expected to fire (see
  \S\ref{sec:severity-collapse}). \\[2pt]
$\rho$-band &
  Monte-Carlo \textsc{go} under $\rho=0$ but \textsc{no-go} under $\rho=0.3$ &
  Declare the gate \emph{fragile}; default to the binary fallback rather than
  claim a \textsc{go} the correlation could flip (\S\ref{sec:rho-band}). \\
\bottomrule
\end{tabular}
\end{table}

The structural property we want from Table~\ref{tab:kill-crosswalk} is that
\emph{no branch is a project-killing null}. The headline confound-attribution
protocol is designed to be independently publishable: a confounding signal detected
as a transportable, power-conditioned protocol stands regardless of downstream
outcomes. Below it, the guarded per-AU $\kappa$ check is a runnable protocol for
measuring whether a frozen vision-language model can weak-label feline FGS action
units at human-rater agreement (the result is pending an independent per-AU vet
anchor, and is only a weak-labeling-capability measurement if that anchor is
independent and the VLM's rubric differs from the vet's --- otherwise it measures
rubric-following); it is retained as kill-tree insurance, becoming the
sole-survivor headline only if the confound leg degrades. The binary-plus-wrapper
spine is the floor that stands today --- the likely v1 ship --- with the graded
layer as upside conditional on the Gate-1-B $\kappa$ pilot passing.

%====================================================================
\section{Power-Calculation-Conditioned Framing of Abstention}
\label{sec:power-conditioned}

The abstention curve (Component~5) is the component most exposed to the project's
single binding constraint: every quantitative claim ultimately draws on the same
modest vet-confirmed anchor (a Gate-0 placeholder of $120$ confirmed images,
$\geq 50$ pain-positive, echoed from \texttt{configs/power.yaml} and to be replaced
before any quantitative claim; this integer is an unconfirmed placeholder over
zero data, not a power-backed figure). Gate~0 runs three zero-data power
calculations: (a) the number of faces needed for a per-AU $\kappa$ CI half-width
$\leq 0.15$ (the committed budget meets this floor); (b) the number needed for a
Learn-Then-Test / conformal NPV $\geq 0.90$ at an abstention rate $\leq 40\%$
\citep{angelopoulos2021ltt}; and (c) whether the $0.39$-threshold operating-point
CI is reportable at the budgeted $n$.

\textbf{The committed Gate-0 verdict on (c) is negative.} At the budgeted anchor the
$0.39$ per-class sensitivity/specificity operating point rests on only $n_{\text{pos}}
= 16$ positives, giving a Clopper--Pearson half-width of $\approx 0.25$; the power
calculation accordingly flags it \texttt{reportable:false}. We therefore do
\emph{not} treat the $0.39$ point as a result of this work: it is carried as an
exploratory, future-work quantity and named as a limitation, not as an open question
to be settled at the present budget.

The Gate-0(b) NPV sizing is, at the present budget, a \emph{placeholder-illustrative}
calculation, not a settled certification. Its inputs are exploratory --- a
placeholder \texttt{vet\_budget} of $120$, a worked total of $84$ faces, and a
target NPV of $0.90$ over zero measured data --- so although the calculation marks
the band as nominally meeting the target, we do \emph{not} report that band as
``certified'' or ``guaranteed.'' The band is not yet certified because the budget is
an unconfirmed placeholder over an empty corpus, not because any arithmetic forbids
it; whether a real, non-placeholder budget can certify a useful NPV floor on
measured data is future work. The decision rule we pre-register is one of
\emph{reporting discipline}: the word ``guaranteed'' never appears anywhere in the
paper, and the abstention curve is reported only as ``NPV with a one-sided $95\%$
lower bound at each abstention rate'' --- a finite-sample statement. The $0.95$
arithmetic is retained for one purpose only: to make the ban concrete. Certifying
NPV~$\geq 0.95$ would, at this scale, demand on the order of $60$ zero-error
abstained-in negatives out of ${\sim}70$ total, which is exactly the kind of
fragile, near-vacuous claim a point ``guarantee'' would smuggle in --- hence the
one-sided lower bound and the placeholder framing rather than a deliverable
certificate. Should the placeholder sizing not survive a real budget, the headline
weight of the wrapper relocates onto the welfare-asymmetric decision curve
(Component~4), the more robust artifact at small $n$; this is a relocation, not a
kill.

%====================================================================
\section{Expected Severity-Cell Collapse (Gate 6)}
\label{sec:severity-collapse}

Gate~6 is a pure-counting gate applied after the vet anchor exists: it counts the
number of vet-confirmed AU=2 (severe) cells per action unit. The pre-registered
rule is that if any AU's $=2$ cell count is single-digit, the high end of the
scale is \emph{collapsed} (merge the $1$ and $2$ levels, or report only
painful/not-above-threshold). This is a structural decision, not a caveat: a
caveat would leave an anchoring number in a table, whereas an AU=2 sensitivity
estimate at single-digit $n$ carries a confidence interval spanning roughly
$[0.35, 0.97]$ --- that is no information.

\textbf{We expect Gate~6 to fire.} Given the ${\sim}12$--$13\%$ pain prevalence of
the source corpus and the further rarity of \emph{severe} grimace within the pain
class, the budgeted anchor is very unlikely to deliver double-digit AU=2 cells.
Planning for the collapse in advance --- rather than treating it as a surprise ---
is what keeps the graded layer honest: no per-cell severe calibration number is
printed, and the graded layer remains an inspected-not-validated artifact. This is
the operational reason the word ``graded'' is struck from every validated-claim
sentence; the low-prevalence measurement-error dynamics that make such tail cells
untrustworthy are exactly those characterised by \citet{chavoshi2025}, whose
analysis shows how a high apparent agreement can coexist with a badly degraded
operating sensitivity once prevalence is low.

%====================================================================
\section{Monte-Carlo Correlation-Sensitivity Band}
\label{sec:rho-band}

A naive plan would estimate a full $5\times5$ cross-AU error-correlation matrix
from the anchor and feed it into the error-propagation Monte-Carlo. At the
budgeted $n$ this is false rigor: a ten-off-diagonal matrix has per-entry
confidence intervals on the order of $[-0.4, +0.6]$, may not even be
positive-definite, and launders sampling noise as data-driven structure. We
therefore replace it with a \textbf{two-point $\rho$-sensitivity band}: the
error-propagation Monte-Carlo is run twice, once under independence ($\rho = 0$)
and once under a pinned uniform correlation ($\rho = 0.3$), and the gate returns
\textsc{go} \emph{only if the decision holds under both}.

The decision rule is pre-registered (Table~\ref{tab:kill-crosswalk}, final row):
if the Monte-Carlo is \textsc{go} under $\rho = 0$ but \textsc{no-go} under
$\rho = 0.3$, we declare the gate \emph{fragile}, say so explicitly, and default to
the binary fallback rather than claim a \textsc{go} that a plausible correlation
could flip. A pinned $\rho$ is an honest sensitivity analysis; a noisily-fit
matrix would be a false claim of empirical precision. The two-point band is the
conservative substitute that we can defend at this sample size; the full matrix is
explicitly deferred (re-added only when $n$ reaches the low hundreds of vet rows
with non-single-digit cells).

%====================================================================
\section{Literature Anchors as Background, Never as Our Results}
\label{sec:lit-anchors}

The numbers in Table~\ref{tab:lit-anchors} are the prior literature's, cited as
\emph{background anchors} that motivate our design choices --- chiefly the
fixed-recall operating point and the anti-benchmark reporting discipline. \textbf{None
of these figures is a result of this work}, and the paper never makes a
``we beat $77/79/95\%$'' claim. We report against the distinct-pain-cat
denominator with Clopper--Pearson / bootstrap confidence intervals, and the
literature figures appear only as prior-art context, never as a head-to-head
comparison.

\begin{table}[t]
\centering
\caption[Literature background anchors]{Literature background anchors. Every figure
below belongs to the cited authors and is reproduced here as background context for
our design --- \emph{not} as a result of this work. We do not benchmark against
these numbers; they motivate the fixed pain-recall $\geq 0.90$ operating point and
the decision-support framing.}
\label{tab:lit-anchors}
\footnotesize
\begin{tabular}{@{}p{0.20\linewidth}p{0.34\linewidth}p{0.36\linewidth}@{}}
\toprule
Source (theirs) & Their reported figure & How we use it (background only) \\
\midrule
\citet{evangelista2019} &
  FGS validation: \evangelistanums{}; pain at a $>0.39$ ($\approx 4/10$) threshold &
  Sources our fixed pain-recall $\geq 0.90$ \emph{target} and the $0.39$ decision
  cutoff; never a comparison baseline. \\[2pt]
\citet{steagall2023} &
  Automated FGS via landmark$\to$geometry$\to$XGBoost, $95.5\%$ accuracy &
  Names a method family (landmark$+$XGBoost) we deliberately do \emph{not}
  replicate; cited to position our conceded-engine choice. \\[2pt]
\citet{feighelstein2023} &
  Binary cat-pain: landmark $77\%$ vs.\ deep-learning $65\%$; XAI mouth $>$ eyes
  $>$ ears &
  Background for the binary spine and for AU saliency expectations; not a target
  to beat. \\[2pt]
\citet{chavoshi2025} &
  Low-prevalence measurement-error propagation (e.g.\ ${\sim}10\%$ prevalence
  $\to$ ${\sim}53\%$ effective sensitivity) &
  Motivates the severity-cell collapse (\S\ref{sec:severity-collapse}) and the
  low-prevalence CI discipline. \\
\bottomrule
\end{tabular}
\end{table}

\paragraph{Restatement of the honesty contract.}
To close the chapter, we restate its governing constraint without hedging.
\emph{No experiments have been run.} This chapter contains no measured results.
The gate-conditioned outcomes above are pre-registered decision rules and the
products that ship on each branch, not observations. The only numbers presented
are the literature anchors of Table~\ref{tab:lit-anchors}, each attributed to its
authors. When the project executes, every reported figure will carry its
distinct-pain-cat denominator and a Clopper--Pearson or bootstrap confidence
interval, sensitivity and specificity at the $0.39$ cutoff --- a quantity the
Gate-0 power calc flags \texttt{reportable:false} at the budgeted
$n_{\text{pos}} = 16$, hence exploratory pending a future enlarged vet anchor ---
will be estimated on vet-confirmed labels only (the circularity firewall), and the
agreement between the model and the weak-labeling VLM will never be presented as
validation. The
modal product --- the binary-plus-wrapper instrument headlined by the
confound-attribution protocol, with the guarded $\kappa$ check below it and an
inspected-not-validated graded layer --- is the outcome we plan for, and it
stands even if the graded layer is dropped entirely.
